\documentclass[aspectratio=169]{beamer}

\usetheme{Madrid}
\usecolortheme{seahorse}

\usepackage{graphicx}
\usepackage{xcolor}
\usepackage{hyperref}

\title{Meeting 2: How the Real System Works}
\subtitle{Autonomy Onboarding}
\author{MIT Arcturus}
\date{}

\begin{document}

\begin{frame}
    \titlepage
\end{frame}

% ============================================================
\section{Big Picture}
% ============================================================

\begin{frame}{Last Time}
    You built a simplified autonomy pipeline:
    \vfill
    \begin{center}
        Known buoy positions $\rightarrow$ Compute midpoint $\rightarrow$ PID control $\rightarrow$ Move boat
    \end{center}
    \vfill
    The real boat does the same thing, but each step is more complex.
    \vfill
    Today we'll walk through the real system and see how it all connects.
\end{frame}

\begin{frame}{The Real Pipeline}
    \begin{center}
        \Large
        Perception $\rightarrow$ Mapping $\rightarrow$ Planning $\rightarrow$ Control
    \end{center}
    \vfill
    \begin{enumerate}
        \item Perception --- cameras and LiDAR figure out what's around us
        \item Mapping --- track those detections over time, build a map
        \item Planning --- decide where to go and find a path there
        \item Control --- follow that path with the motors
    \end{enumerate}
    \vfill
    A task manager sits on top and tells the system what to do (``go through the gates,'' ``circle the buoy,'' etc.)
\end{frame}

% ============================================================
\section{Perception}
% ============================================================

\begin{frame}{Perception: What's Around Us?}
    The boat has two main sensors:
    \vfill
    \begin{itemize}
        \item Cameras --- see color and shape (like your eyes)
        \item LiDAR --- measures distance to objects using lasers (like sonar but with light)
    \end{itemize}
    \vfill
    Cameras are good at recognizing what something is (red buoy vs green buoy).\\
    LiDAR is good at knowing exactly where something is in 3D.\\
    \vfill
    We combine both to get labeled 3D positions of obstacles.
\end{frame}

\begin{frame}{Step 1: Object Detection (YOLO)}
    A neural network called YOLO looks at each camera frame and draws bounding boxes around objects it recognizes.
    \vfill
    Output: ``there's a red buoy in this rectangle of the image''
    \vfill
    This tells us what and roughly where in the image, but not how far away.
    \vfill
    % TODO: screenshot of YOLO detections on a camera frame
    \begin{center}
        \textit{(YOLO detection screenshot here)}
    \end{center}
\end{frame}

\begin{frame}{Step 2: Depth from LiDAR}
    LiDAR gives us a 3D point cloud --- thousands of distance measurements.
    \vfill
    We project the LiDAR points onto the camera image. Points that land inside a YOLO bounding box get that box's label.
    \vfill
    Now we know: ``these 3D points are a red buoy.''
    \vfill
    This happens in the \texttt{bbox\_project\_pcloud} node.
    \vfill
    % TODO: diagram showing camera image + lidar points overlay
    \begin{center}
        \textit{(LiDAR + camera fusion diagram here)}
    \end{center}
\end{frame}

\begin{frame}{Step 3: Obstacle Tracking}
    Raw detections are noisy --- objects appear and disappear between frames.
    \vfill
    The \texttt{object\_tracking\_map} node tracks obstacles over time:
    \begin{itemize}
        \item Matches new detections to previously seen obstacles
        \item Smooths positions across frames
        \item Publishes an \texttt{ObstacleMap} message with all tracked obstacles
    \end{itemize}
    \vfill
    Each obstacle has: a label (what it is), a position (where it is), and a convex hull (its shape).
\end{frame}

\begin{frame}{Perception Summary}
    \begin{center}
        Camera frame $\rightarrow$ YOLO (labels) $\rightarrow$ LiDAR fusion (3D positions) $\rightarrow$ Tracking (stable map)
    \end{center}
    \vfill
    The output is an \texttt{ObstacleMap} --- a list of labeled obstacles with positions.
    \vfill
    This is what your code subscribes to when writing autonomy tasks. You don't need to worry about how perception works internally, just know that you'll get a list of obstacles with labels and positions.
\end{frame}

% ============================================================
\section{Localization and Mapping}
% ============================================================

\begin{frame}{Where Are We?}
    Perception tells us what's around us, but we also need to know where we are in the world.
    \vfill
    Two sources:
    \begin{itemize}
        \item GPS --- gives a global position (latitude/longitude, converted to x/y)
        \item SLAM --- Simultaneous Localization and Mapping, refines our position by matching sensor data over time
    \end{itemize}
    \vfill
    Together they give us the boat's position and heading in a global coordinate frame.
\end{frame}

\begin{frame}{Coordinate Frames}
    The system uses two main frames:
    \vfill
    \begin{itemize}
        \item \texttt{map} --- the fixed world frame (doesn't move)
        \item \texttt{base\_link} --- attached to the boat (moves with it)
    \end{itemize}
    \vfill
    Perception gives obstacles in \texttt{base\_link} (relative to the boat). \\
    Using localization, we convert those into \texttt{map} coordinates (global positions).
    \vfill
    The obstacle map has both: \texttt{local\_point} (in base\_link) and \texttt{global\_point} (in map).
\end{frame}

\begin{frame}{Global Map}
    The \texttt{object\_tracking\_map} node combines:
    \begin{itemize}
        \item Local obstacle detections (from perception)
        \item Robot position (from localization)
    \end{itemize}
    \vfill
    To produce a global obstacle map --- obstacle positions in the world frame that persist as the boat moves.
    \vfill
    There's also a grid map (\texttt{grid\_map\_generator}) that creates an occupancy grid for path planning --- think of it as a bird's-eye view where cells are either free or occupied.
\end{frame}

% ============================================================
\section{Planning and Control}
% ============================================================

\begin{frame}{Planning: Where Do We Go?}
    Given a goal position (e.g., midpoint between two buoys) and a map of obstacles, the planner finds a path.
    \vfill
    The \texttt{navigation\_server} uses A* search on the occupancy grid:
    \begin{itemize}
        \item Finds a path from the boat's position to the goal
        \item Avoids obstacles on the grid
        \item Returns a list of waypoints to follow
    \end{itemize}
    \vfill
    In your lab, you drove directly to the midpoint with PID. On the real boat, A* finds a path around obstacles first.
\end{frame}

\begin{frame}{Control: How Do We Get There?}
    The \texttt{controller\_server} takes the planned path and drives the boat along it.
    \vfill
    \begin{itemize}
        \item PID control (like you wrote) to follow each waypoint
        \item Potential field obstacle avoidance --- if something unexpected is close, steer away from it in real time
        \item Outputs velocity commands to the thrusters
    \end{itemize}
    \vfill
    This is the same idea as your lab (PID to a point), just with extra safety layers.
\end{frame}

\begin{frame}{Planning + Control Summary}
    \begin{center}
        Goal position $\rightarrow$ A* path planning $\rightarrow$ Path following (PID) $\rightarrow$ Thruster commands
    \end{center}
    \vfill
    In your lab:
    \begin{center}
        Goal position $\rightarrow$ PID directly $\rightarrow$ Velocity commands
    \end{center}
    \vfill
    The real system adds a planner between the goal and the PID to avoid obstacles along the way.
\end{frame}

% ============================================================
\section{Task Execution}
% ============================================================

\begin{frame}{Tasks: What Should We Be Doing?}
    Everything so far answers ``how do we get from A to B.''
    \vfill
    Tasks answer ``what is A and B in the first place?''
    \vfill
    Each competition task (entry gates, follow the path, docking, etc.) is a separate program that:
    \begin{enumerate}
        \item Reads the obstacle map
        \item Decides what to do (find the next buoy pair, pick a dock, etc.)
        \item Sends goal positions to the planner
    \end{enumerate}
\end{frame}

\begin{frame}{Task Structure}
    Each task is a Python class that inherits from \texttt{TaskServerBase}.
    \vfill
    Two methods to implement:
    \begin{itemize}
        \item \texttt{init\_setup()} --- runs once at the start (find initial buoys, etc.)
        \item \texttt{control\_loop()} --- runs repeatedly (track buoys, send goals, check if done)
    \end{itemize}
    \vfill
    \texttt{TaskServerBase} gives you helpers like:
    \begin{itemize}
        \item \texttt{move\_to\_point()} --- send the boat to a position
        \item \texttt{get\_robot\_pose()} --- get the boat's current position
        \item \texttt{mark\_successful()} --- signal that the task is done
    \end{itemize}
\end{frame}

\begin{frame}{Task Manager}
    The \texttt{run\_tasks} node is the state machine that runs everything.
    \vfill
    \begin{itemize}
        \item Decides which task to execute next
        \item Sends start/stop signals to task servers
        \item Handles interrupts (e.g., harbor alert --- drop everything and respond to a signal)
    \end{itemize}
    \vfill
    Each task runs as an action server. The task manager is the action client that tells them when to start.
\end{frame}

\begin{frame}{The Full Picture}
    \begin{center}
        \small
        Task Manager $\rightarrow$ Task Server (e.g., Follow Path)
    \end{center}
    \begin{center}
        $\downarrow$
    \end{center}
    \begin{center}
        \small
        ``Go to this point'' $\rightarrow$ Navigation Server (A* planning)
    \end{center}
    \begin{center}
        $\downarrow$
    \end{center}
    \begin{center}
        \small
        Planned path $\rightarrow$ Controller Server (PID + avoidance) $\rightarrow$ Thrusters
    \end{center}
    \vfill
    Meanwhile, perception is continuously running in the background, updating the obstacle map that everything else reads from.
\end{frame}

% ============================================================
\section{Codebase Tour}
% ============================================================

\begin{frame}{Repository Structure}
    The codebase is split into packages:
    \vfill
    \begin{itemize}
        \item \texttt{all\_seaing\_perception} --- camera, LiDAR, YOLO, point cloud fusion, tracking
        \item \texttt{all\_seaing\_navigation} --- A* planner, grid map
        \item \texttt{all\_seaing\_controller} --- PID, potential field, thruster control
        \item \texttt{all\_seaing\_autonomy} --- competition task implementations
        \item \texttt{all\_seaing\_common} --- base classes (TaskServerBase, ActionServerBase)
        \item \texttt{all\_seaing\_bringup} --- launch files and config
        \item \texttt{all\_seaing\_driver} --- hardware interfaces
    \end{itemize}
\end{frame}

\begin{frame}{Explore the Codebase}
    Try to find these in the \texttt{all\_seaing\_vehicle} repo:
    \vfill
    \begin{enumerate}
        \item Which node runs YOLO object detection?
        \item What message type does the obstacle map use?
        \item Where is \texttt{TaskServerBase} defined? What methods does it provide?
        \item Find the follow-the-path task. What are \texttt{init\_setup()} and \texttt{control\_loop()} doing?
        \item How does \texttt{run\_tasks.py} decide which task to run next?
        \item Where are the PID gains configured?
    \end{enumerate}
    \vfill
    Use \texttt{grep}, file search, or just browse the repo. Ask questions!
\end{frame}

\begin{frame}{Next Time}
    Now you know how all the pieces fit together.
    \vfill
    Next meeting we'll start working with the real perception pipeline --- running it on recorded data and writing a controller that uses actual obstacle detections.
\end{frame}

\end{document}
